\documentclass[10pt,a4paper]{amsart}
\usepackage[margin=19mm]{geometry}
\usepackage{amsmath,amssymb,mathtools}
\usepackage[scr=rsfs,cal=euler]{mathalfa}
\usepackage{xfrac}
\usepackage[unicode,hidelinks,bookmarks=false]{hyperref}
\newcommand{\ie}{i.e.\ }
\newcommand{\cf}{cf.\ }
\newcommand{\resp}{respectively\ }
\newcommand{\Subsection}{\subsection}
\providecommand{\nobreakdash}{\nobreak\hbox{-}}
\setlength{\emergencystretch}{2em}
\multlinegap0pt
\usepackage{bm}
\usepackage{bbm}
\usepackage{dsfont}
\usepackage{xspace}
\usepackage{cancel}
\usepackage{mathtools}
\usepackage{amsthm}
\makeatletter
\@ifundefined{definition}{\newtheorem{definition}{Definition}}{}
\@ifundefined{lemma}{\newtheorem{lemma}{Lemma}}{}
\makeatother
\DeclareMathOperator{\Aff}{Aff}
\newcommand{\C}{\mathbb{C}}
\newcommand{\HH}{\mathbb{H}}
\newcommand{\Nc}{\mathcal{N}}
\newcommand{\OO}{\mathcal{O}}
\newcommand{\R}{\mathbb{R}}
\newcommand{\Z}{\mathbb{Z}}
\title[Automorphism groups of Inoue surfaces $S^{(+)}/S^{(-)}$]{Automorphism groups of Inoue surfaces $S^{(+)}/S^{(-)}$}
\author{David Petcu}
\date{Source: arXiv:2509.06860v1 (CC BY 4.0)}
\begin{document}
\maketitle

\noindent\emph{Source and licence.} Automorphism groups of Inoue surfaces $S^{(+)}/S^{(-)}$. Version arXiv:2509.06860v1 (CC BY 4.0), distributed under the Creative Commons Attribution 4.0 International licence, as recorded in the arXiv licence metadata for this exact version and shown on the abstract page https://arxiv.org/abs/2509.06860v1. Full rights notice: licenses/arxiv-2509.06860-CC-BY-4.0.txt.\par\smallskip

\noindent\textit{Editorial scope.} Counted unit: the Lemma whose source label is \texttt{norm} in the article (source0.tex lines 570--572, source label \texttt{norm}), delivered together with the article's own complete original proof of it (source0.tex lines 574--646, one \texttt{proof} environment of 73 lines). No mathematics is written, repaired or re-derived: statement and proof are the article's own text line for line. Required context retained uncounted: Incl (lines 357--359). Every definition, display and label of those lines is retained unchanged. Omitted from this package and not redistributed: 1--356, 360--569, 573--573, 647--1017. Nothing inside a retained excerpt is omitted or altered: every retained body is delivered exactly as the article's lines are written, so no omission receipt of any kind accompanies this package. Citations: the retained excerpts carry no citation command at all, so no bibliography entry of the article is transcribed and no reference is redistributed. Reference adaptation: none was needed and none was made; every reference inside the delivered unit resolves inside it, so no unresolved reference or citation is produced. Printed slips kept literal: no printed word, symbol, hypothesis or number of the counted statement or of its proof was altered. Layout and class: the article's own class is \texttt{amsart} with packages including geometry, amsmath, amssymb, amsthm, parskip, tikz-cd, mathtools, enumitem, fontenc, lmodern, babel, float, dsfont, natbib; the wrapper is the lane's pinned \texttt{amsart} editorial wrapper at 10pt on A4, which restates the theorem-like environments the retained text needs and adds the article's own macro definitions that the retained text needs (\texttt{\textbackslash Aff}, \texttt{\textbackslash C}, \texttt{\textbackslash HH}, \texttt{\textbackslash Nc}, \texttt{\textbackslash OO}, \texttt{\textbackslash R}, \texttt{\textbackslash Z}), with extra wrapper packages (bm, bbm, dsfont, xspace, cancel, mathtools). The delivered numbering is the unit's own. Assets: no retained span contains a figure environment, a graphics-inclusion command, a PDF-page inclusion, a table or a TikZ picture; no image file is copied into this package, the package declares no asset dependency and no image file is redistributed with it. Licence: version arXiv:2509.06860v1 of this article is distributed under the Creative Commons Attribution 4.0 International licence, as recorded in the arXiv licence metadata for this exact version and shown on the abstract page https://arxiv.org/abs/2509.06860v1. A full rights notice is delivered at licenses/arxiv-2509.06860-CC-BY-4.0.txt. Characters: every retained excerpt is pure ASCII, so the delivered engine, XeLaTeX with its default Latin Modern text font, reports no missing character.
\begin{definition}\label{Incl}
Let $\mathcal{G}_{\theta}^{(+)}$ be the group $\OO_K^{\times, >} \times K \times \C$, with the group structure induced by the above embedding $\iota_{\theta}^{(+)}$ in $\mathcal{G}$.
\end{definition}

\begin{lemma}\label{norm} 
Let  $\Nc(\Gamma)$ be the normalizer of $\Gamma$ in $\Aff(\HH \times \C)$. Then, $\Nc(\Gamma)$ is contained in $\mathcal{G}^{(+)}_\theta$ (\ref{Incl}). Moreover, if $h$ is an element of $\Nc(\Gamma)$, then, viewed as an element of $\mathcal{G}_{\theta}^{(+)}$, $h$ is of the form $[\mu,l,s]$ with $\mu \in {}^I[\OO_K^{\times, >}]$, $l \in I(1-u_\theta)^{-1}$ and $s \in \C$.
\end{lemma}

\begin{proof}
There is a surjective morphism from $\Aff(\HH \times \C)$ to the group
$$G=
\left\{
\left[
\begin{pmatrix}
\mu_1 & 0\\
0 & \mu_2
\end{pmatrix},
\begin{pmatrix}
l_1 \\
l_2
\end{pmatrix}
\right] \in \text{GL}(2, \C) \ltimes \C^2\vert \mu_1\in \R_>;\  \mu_2\in \C^*;\  l_1, l_2 \in \R \right\}$$
sending
$$\left[
\begin{pmatrix}
\mu & 0\\
\lambda & \nu
\end{pmatrix},
\begin{pmatrix}
l \\
s
\end{pmatrix}
\right] \to
\left[
\begin{pmatrix}
\mu & 0\\
0 & \frac{\nu}{\mu}
\end{pmatrix},
\begin{pmatrix}
l \\
\frac{\lambda}{\mu}
\end{pmatrix}
\right]
$$
The kernel is $T_\C$, so $\Aff(\HH \times \C)/T_\C$ is isomorphic to $G$. The group $\Gamma/(\Gamma \cap T_\C)$ is isomorphic to $\langle u_\theta \rangle \ltimes I$. Let $\hat{g_i}$ denote the class of the generator $g_i$ in the quotient $\Gamma / (\Gamma \cap T_\C)$. Now let $h \in \Nc(\Gamma)$ with 
$$\hat{h} =
\left[
\begin{pmatrix}
\mu_1 & 0\\
0 & \mu_2
\end{pmatrix},
\begin{pmatrix}
l_1 \\
l_2
\end{pmatrix}
\right]
$$
the image of $h$ in $G$. Then $\hat{h}$ is contained in the normalizer of $\Gamma/(\Gamma \cap T_\C)$ in $G$. Which, in particular, means that $\hat{h}\hat{g_0}\hat{h}^{-1} \in \Gamma/(\Gamma\cap T_\C)$, so 
$$\left[
\begin{pmatrix}
\sigma_1(u_\theta) & 0\\
0 & \sigma_2(u_\theta)
\end{pmatrix},
\begin{pmatrix}
l_1\sigma_1(1-u_\theta) \\
l_2\sigma_2(1-u_\theta)
\end{pmatrix}
\right] \in \langle u_\theta \rangle \ltimes I
$$
Thus there exists $l \in K$ such that $\sigma_i(l)=l_i$, $i \in \{1,2\}$. In fact, $l$ is an element in $I(1-u_\theta)^{-1}$.
Similarly, we have $\hat{h}\hat{g_i}\hat{h}^{-1} \in \Gamma/(\Gamma\cap T_\C)$, so
$$
\left[
I_2,
\begin{pmatrix}
\mu_1\sigma_1(x_i) \\
\mu_2\sigma_2(x_i)
\end{pmatrix}
\right] \in \langle u_\theta\rangle \ltimes I,\text{ for }i\in\{1,2\}$$
hence, once again, there exists some element $\mu \in K$ such that $\sigma_i(\mu)=\mu_i$. Moreover, $\mu$ is a unit in $\OO_K$, $\sigma_1(\mu)>0$ and $I$ a $\Z[\mu]$-fractional ideal, thus $\mu \in {}^I[\OO_K^{\times,>}]$. We have shown that $h = [\mu, l, s]\in \mathcal{G}_{\theta}^{(+)}$, where $\mu \in {}^I[\OO_K^{\times,>}]$, $l \in I(1-u_\theta)^{-1}$ and $s$ is some complex number.
\end{proof}

\end{document}
